\section{Background}\label{sec:background}

\subsection{Two forms of Okun's law}

\citet{okun1962} related the unemployment rate to the gap between actual and potential
output. The \emph{difference} form used here avoids estimating potential output. For
country $i$ and quarter $t$, let $u_{it}$ be the unemployment rate (\% of the labour force),
$\Delta u_{it}=u_{it}-u_{i,t-1}$ its change in percentage points (pp), and $g_{it}$ the
quarter-on-quarter growth of real GDP (\%). Then
\begin{equation}\label{eq:okun}
  \Delta u_{it} = \alpha_i + \sum_{k=0}^{K}\beta_{ik}\, g_{i,t-k} + \varepsilon_{it},
  \qquad \theta_i \equiv \sum_{k=0}^{K}\beta_{ik}.
\end{equation}
The coefficient sum $\theta_i$ is the change in the unemployment rate, in pp, by the end of
quarter $t+K$ that goes with one extra point of growth in quarter $t$ alone, which leaves
the level of output about one per cent higher. Under a sustained one-point rise in quarterly
growth, $\theta_i$ is instead the change in $\Delta u$ per quarter once all lags have
passed, and the cumulative change in the rate keeps growing. $-\alpha_i/\theta_i$ is the
growth rate
consistent with stable unemployment. Lagged growth allows the labour market to respond with
a delay. Estimated coefficients are negative in advanced economies, but their size
differs widely from country to country \citep{ball2017okun}.

\subsection{An accounting identity}

Output $Y$ is the product of output per hour, hours per worker and employment. Employment
$E$ is the labour force $L$ times one minus the unemployment rate. In log changes (lower
case) and with the rate approximation,
\begin{equation}\label{eq:identity}
  g \;\approx\; \underbrace{\Delta\ln(Y/H)}_{\text{productivity}}
  + \underbrace{\Delta\ln(H/E)}_{\text{hours per worker}}
  + \underbrace{\Delta\ln E}_{\text{employment}},
  \qquad
  \Delta u \;\approx\; -(\Delta\ln E - \Delta\ln L).
\end{equation}
An Okun coefficient of $-1$ would require productivity, hours per worker and the labour
force not to move with output. Observed coefficients are much smaller in magnitude because
output fluctuations are absorbed through other margins:
\begin{itemize}
  \item productivity (labour hoarding);
  \item hours per worker (short-time work);
  \item participation, through discouraged or added workers.
\end{itemize}
Institutions shape these margins. Strict dismissal rules combined with flexible temporary
contracts can make employment more volatile \citep{bentolila2012twotier}. Unemployment
inflows and outflows differ widely across countries \citep{blanchard2001portugal}, and
hours per worker vary along the cycle \citep{ohanian2012hours}. The identity also shows
that measurement matters: GDP is a domestic concept (production on the territory), whereas
unemployment and participation describe residents.

\subsection{What a reduced-form coefficient measures}\label{sec:weighted}

Let the reduced-form innovations of $(g,\Delta u)$ be driven by mutually uncorrelated
structural shocks $\varepsilon_k$ (demand, productivity, labour supply, energy, \dots):
\[
  g=\sum_k a_k\varepsilon_k, \qquad \Delta u=\sum_k b_k\varepsilon_k .
\]
Then the regression coefficient of $\Delta u$ on $g$ is
\begin{equation}\label{eq:weighted}
  \beta_{\mathrm{Okun}}=\frac{\operatorname{Cov}(\Delta u,g)}{\operatorname{Var}(g)}
  =\sum_k w_k\,\kappa_k,\qquad
  \kappa_k=\frac{b_k}{a_k},\quad w_k=\frac{a_k^2}{\sum_j a_j^2}.
\end{equation}
The coefficient is a weighted average of shock-conditional Okun coefficients
$\kappa_k$. The weights are each shock's share of output variance, and the same algebra holds
horizon by horizon for cumulative responses. Differences between countries or periods can
therefore reflect different labour-market responses ($\kappa_k$) or different shock
mixes ($w_k$). The data cannot separate the two without identifying the shocks, which is
the subject of section~\ref{sec:identification}. Until then, every coefficient is a
\emph{descriptive association}.
